% ---
% id: meat-core/cor:reduction-constrained
% kind: cor
% title: "The reduction for arbitrary constraint lists"
% summary: >-
%   For any declared constraint list, licensed or not, the constrained value
%   bounds the constrained single-round rate from below, with the same list on
%   both sides. A key length from the constrained value still needs the
%   attack-set licence of [[thm:constrained-accumulation]].
% prominence: standard
% status: implemented
% uses: [def:homog, def:hqkd-list, def:feasible, thm:reduction, def:admissible-declaration]
% mentions: [thm:constrained-accumulation]
% assumes: []
% refs: []
% provenance: newly-derived
% implements: certiqs_key.meat_core.objective.objective:constraint_value
% lean:
% ---

No step of the reduction reads the constraint list beyond checking that the
joint dephasing keeps every entry, so the comparison holds for any list.
Whether the constrained value may enter a key length is a separate question,
answered by the attack-set licence (Theorem~\ref{thm:constrained-accumulation}).

\begin{corollary}\label{cor:reduction-constrained}
Let $\CList$ be arbitrary: $\Theta_k$ an arbitrary Hermitian operator on $A$
in \tref{trusted-source mode}, $\Theta_k=c_k$ an arbitrary real number in
\tref{untrusted-source mode}, and entries declared \tref{party-local} in
either mode (Definition~\ref{def:admissible-declaration}\,(vi)); no licensing restriction
(Definition~\ref{def:homog}) is needed in either mode. Then
\begin{equation}\label{eq:c8-general}
  \hqkd(\CList)\;\ge\;\hlow(\CList),
\end{equation}
where $\hqkd(\CList)$ is the \tref{constrained single-round rate} of
Definition~\ref{def:hqkd-list} and $\hlow(\CList)$ denotes
\eqref{eq:hlow-def} over the feasible set $\Feas$ of \eqref{eq:feas},
which carries the same list ($\hlow(\CList)$ and the unadorned
$\hlow$ are the same quantity; the argument is written when the
dependence on the declared list matters, as in
$\hlow(\CList)\ge\hlow(\varnothing)$).
\end{corollary}

\begin{proof}
Steps 1--5 of the proof of Theorem~\ref{thm:reduction} apply verbatim
to a feasible point of \eqref{eq:hqkd-list}: Step 1 preserves every
entry. For an entry read through its lifted operator, block diagonality
of $\hat O_k$ along $\SecSet$ and self-adjointness of $\PinS$ give
$\tr_{\Qsys}[(\idon{A}\otimes\hat O_k)\,\omega'_{A\Qsys}]
=\mathcal{P}_{\mathbb{T}}\big(\tr_{\Qsys}[(\idon{A}\otimes\hat O_k)\,
\omega_{A\Qsys}]\big)$; in trusted-source mode $\mathcal{P}_{\mathbb{T}}$ is
positive and linear and fixes $\Theta_k$, which is block diagonal along
$\mathbb{T}$ ($c_k\tau_A$ for a \tref{homogeneous} entry, by type for
a \tref{general} one, Definition~\ref{def:admissible-declaration}\,(iii)), so
$X\bowtie_k\Theta_k$ implies
$\mathcal{P}_{\mathbb{T}}(X)\bowtie_k\Theta_k$; in untrusted-source mode the
entry value is the trace of the same operator, which the pinching
preserves. For a party-local entry scoped to $f$, with
$\hat D=\hat O_k-c_k\idon{\Qsys_f}$ block diagonal for the pinching
$\mathcal{P}_f$ of that factor and $\mathcal{P}_{f^c}$ the product of the
other factors' pinchings,
$\tr_{\Qsys_f}[(\idon{}\otimes\hat D)\,\omega'_{A\Qsys}]
=(\mathcal{P}_{\mathbb{T}}\otimes\mathcal{P}_{f^c})\big(
\tr_{\Qsys_f}[(\idon{}\otimes\hat D)\,\omega_{A\Qsys}]\big)$, and a positive
map preserves $\bowtie_k0$; Steps 2--3 do not change
the $A\Qsys$-marginal, and the state
$\tilde\omega=\omega'_{A\Qsys}$ of Step 4 therefore satisfies every
entry of $\CList$ in addition to $\tr\tilde\omega=1$, block diagonality along the joint
sectors
and (in trusted-source mode) $\tilde\omega_A=\tau_A$, that is,
$\tilde\omega\in\Feas$. Taking infima over
the constrained sets on both sides yields \eqref{eq:c8-general}.
\end{proof}
